\documentclass[11pt,letterpaper]{article}
\usepackage[T1]{fontenc}
\usepackage[utf8]{inputenc}
\usepackage{lmodern,microtype}
\usepackage[margin=.92in]{geometry}
\usepackage{amsmath,amssymb,amsthm,mathtools}
\usepackage{booktabs,longtable,array,tabularx,enumitem}
\usepackage{xcolor,xurl,hyperref,fancyhdr,needspace}
\definecolor{ink}{HTML}{16364B}
\definecolor{teal}{HTML}{126B76}
\hypersetup{colorlinks=true,linkcolor=ink,urlcolor=teal,citecolor=ink,
 pdftitle={PA194QNS: construction, complete corpus comparison, and reduction research},
 pdfauthor={Research prepared for Vladimir Reshetnikov}}
\setlength{\emergencystretch}{3em}
\setlength{\parindent}{0pt}
\setlength{\parskip}{4pt}
\setlist{itemsep=3pt,topsep=5pt,leftmargin=1.6em}
\pagestyle{fancy}\fancyhf{}
\fancyhead[L]{\small PA consistency machines}
\fancyhead[R]{\small Construction and exact-table evidence}
\fancyfoot[C]{\thepage}
\setlength{\headheight}{14pt}
\newcommand{\file}[1]{\nolinkurl{#1}}
\newcommand{\PA}{\mathsf{PA}}
\newcommand{\Con}{\operatorname{Con}}
\newcommand{\Halt}{\operatorname{Halt}}
\newcommand{\Prf}{\operatorname{Prf}}
\newtheorem{theorem}{Theorem}[section]
\newtheorem{lemma}[theorem]{Lemma}
\newtheorem{proposition}[theorem]{Proposition}
% Exact data are generated by build_report.py.

\title{\color{ink}A new 194-state PA-consistency machine\\[5pt]
\large PA194QNS, a complete comparison of the current corpus,\\
and structural reduction research}
\author{Research prepared for Vladimir Reshetnikov}
\date{13 November 2026, Pacific time}
\begin{document}
\maketitle
\begin{abstract}
No verified machine below 194 working states was obtained. The requested
same-size alternative is supplied by \textbf{PA194QNS}, a deterministic
binary machine on one two-way tape, started on the all-zero tape, with
194 working controls and 388 transitions. It combines the Q axiom emitter,
no-push modus ponens, and the symmetric implementation of the original
sum-based Shift test. These mechanisms are attributed to their earlier
packets; their exact composition, source equivalence, transition table,
novelty comparison and fresh verification are the new work.

The source has 3,651 primitive instructions. Its first 3,649 instructions
are literally identical to the pinned QN source, and the original
five-instruction tail is replaced by two instructions. A simple invariant
and symmetry of destructive equality give a PA-formalizable finite-trace
equivalence. The physical count is $135+67-8=194$. All 4,751 arbitrary-register
source outcomes pass both symbolic checks; the independent literal checker
rederives its four list-sweep lemmas from the actual rows. Exact blank
initialization takes 58,004 steps and is checked on the complete tape.

The resulting ordinary mathematical argument, composed with the attributed
QN calibration, gives
$\PA\vdash(\neg\Halt(\mathrm{PA194QNS},\mathrm{blank})\leftrightarrow
\Con(\PA))$. This is not a newly compiled end-to-end Lean or Rocq theorem,
and no completed iPRA consistency reduction is claimed. The comparison
covers all 54 current deliveries and 39 existing 194-state graph classes,
including auxiliary and rejected graphs. The archive supplies the complete
machine, reproducible code, certificates, source proofs, search evidence and
this report. A separate mathematical coding experiment is explicitly
distinguished from the accepted 194-state result.
\end{abstract}
\clearpage
\tableofcontents
\clearpage

\section{Result and the meaning of the claim}

The principal table is \file{construction/machine/machine.tm}. Its working
states are $0,\ldots,193$, its initial state is $0$, and $194$ is the
separate terminal halt label. The head starts at coordinate zero on an
all-zero two-way infinite binary tape. Every action writes one bit and
moves exactly one square left or right. No stationary moves, nonblank
input or change of counting convention is used.

\begin{theorem}[The delivered same-size alternative]
For the exact table in Appendix~\ref{app:table},
\[
 \boxed{\PA\vdash\left(
 \neg\Halt(\mathrm{PA194QNS},\mathrm{blank})\ 
 \longleftrightarrow\ \Con(\PA)\right).}
\]
\end{theorem}

Here $\Halt$ is the ordinary arithmetization of a finite literal-machine
run, and $\Con(\PA)$ abbreviates the absence of a finite PA proof of $0=1$
in a fixed standard proof calculus. The new part of the argument is an
explicit finite-trace transport to the pinned QN implementation; its
ordinary PA consistency calibration is an attributed dependency.
Section~\ref{sec:logical-background} states the logical and proof-word
interfaces rather than inferring them from test execution.

The biconditional is a statement about what PA proves. To infer actual
nonhalting in the standard model one additionally assumes PA's consistency.
The construction does not prove PA consistent, and it does not establish
that 194 is minimal. The unsuccessful searches below have precisely stated
finite or bounded scopes.

The new machine differs from every existing 194-state literal graph even
after arbitrary state renaming, change of spanning root, left/right
reflection and global symbol interchange. More substantively, it changes
a recurring comparator computation relative to QN, and retains a different
quantifier-distribution head from the old sum-NS implementation. It is a
new combination of known mathematical and implementation mechanisms,
not a claim to have invented those mechanisms individually.

\section{What the complete comparison establishes}
\label{sec:comparison}

The comparison is frozen at ProveIt commit
\begin{quote}\small\nolinkurl{317b07268681b44ff352d173a162e2a8c241b2fe}.\end{quote}
It covers all 51 delivered packets filed under
\file{Computability/BusyBeaver/Research/PA/docs}, the three current incoming
archives, the primary PA documentation and imported PA196/PA209 sources.
The ZFC and ConciseIPRA projects were examined for transferable mechanisms.
The three additional archives are \file{PA194NS_research.zip},
\file{pa194haz_research_2026-11-14.zip}, and
\file{pa194qn_research_2026-11-13.zip}. Thus there are \textbf{54 deliveries}
in the narrative ledger. Immutable paths, packet summaries and provenance
appear in \file{comparison/packet_ledger.csv} and its JSON counterpart.

\subsection{Literal tables, graph classes and evidence status}

The new scanner reads every tracked \file{.tm}, \file{.txt}, \file{.csv}
and \file{.tsv} file in the two requested locations, and recursively reads
ZIP archives. It visits identical ZIP payloads once. A recognized table is
a complete numeric five-column binary transition table with left/right
moves and one undefined target serving as halt. The resulting inventory is:
\begin{center}
\begin{tabular}{@{}lr@{}}\toprule
Item & Count\\\midrule
Delivery packets & 54\\
Recognized literal-table occurrences & 1,402\\
Distinct table byte strings & 721\\
Complete action graphs with a spanning root & 497\\
194-state action graphs & 39\\
195-state action graphs & 19\\
196-state action graphs & 99\\
\bottomrule\end{tabular}
\end{center}
The scanner examines 3,480 candidate text occurrences and 17 distinct ZIP
bodies reached through 29 references. The 194-state portion contains 179
occurrences and 94 byte presentations. Two disconnected synthetic fixtures
with six and nine controls are recorded separately; another synthetic
fixture has no terminal target and fails the recognizer. There is
\emph{no complete spanning-root graph with 7--193 working states} in the
recognized corpus. This is an inventory result, not a minimality theorem.

The graph comparison is stronger than a comparison of filenames or hashes.
For every possible root reaching the whole table it serializes the complete
action graph in deterministic traversal order. It then compares the
canonical strings, also allowing global tape reflection and interchange of
the two symbols. All 39 current 194-state classes remain distinct under
these transformations. Symbol interchange is included only as a stronger
structural exclusion test: it does not by itself preserve a distinguished
zero blank. The new machine is compared against \emph{all} these classes,
including auxiliary and rejected objects.

Of the 39 existing classes, \textbf{26 are presented as constructions},
nine are auxiliary or search objects, and four are explicitly rejected
implementations. These labels report the supplied evidence; they do not
assert that all historical proof obligations were independently discharged
in this investigation. Appendix~\ref{app:frontier} lists each class. Every
table, its identity, its member paths and its evidence classification can
also be inspected in the machine-readable census.

Old reports have different occurrence totals because repository filing has
removed repeated table copies, and older archives recursively quoted one
another. Those totals must not be added. The current scanner deliberately
does not infer tables hidden only in arbitrary program expressions, JSON,
prose or uninstantiated search parameters. The packet-by-packet narrative
review complements that explicit literal-table boundary.

\subsection{The main 194-state mechanisms}

\begin{longtable}{@{}p{.23\linewidth}p{.71\linewidth}@{}}
\toprule Family & What changes, and what its comparison requires\\\midrule\endhead
PA194, AD, 68/8 & The separated diagonal head and ten-entry AX interface are shared. The primary budget is $134+67-7$; AD changes register placement, and 68/8 trades an extra decoder reader for an extra compatible fold. Their resolved sources or source conjugacies must be compared, not only the counts.\\
C, CS & Successful MP leaves its result in the cache; Shift tests the cache alone. CS changes the destructive comparator. These have 3,649 and 3,646 source sites. Their proof-word interface differs from the original sum-guard controller.\\
N, H, HS & N omits a successful-MP zero push but retains the sum guard and visible result. H/HS combine no-push MP with a cache-only guard. N's successful result is $(b,w,b,0)$; the retained offset and missing zero must be handled by the proof compiler.\\
M4R, M5, M6; N5, N6 & A multiplier-decoder reassociation exposes another history fold. M5/M6 change formula coefficients to 5/6; the stack and scheduler still use Cantor pairing. Their hereditary formula-and-label recoding is not order preserving. N5/N6 add no-push MP.\\
J, D, Q, QN & J changes return organization; D emits a diagonal instance indexed by another register. Q exchanges the curried premises of a quantifier-distribution axiom. QN adds no-push MP to Q. Logical and command interfaces must both be translated.\\
CZ, CSZ/X, A, F, R, ZR and auxiliaries & Known-zero decrements, scanner choices and source reassociations change physical return paths. Their local invariants need a separate blank-start check. Some source-valid changes fail during initialization.\\
HA, HAZ & These combine the older MPH cached-result interface with physical A/D placement; HAZ also changes the final AX return. Their sum guard differs from the cache-only H/HS family despite overlapping names.\\
SUMS, incoming sum-NS & These retain the original sum guard but compare its two operands in reverse order, removing a transfer loop. A new one-bit history fold pays for the trailer reader. The incoming sum-NS additionally omits the MP push.\\
\bottomrule
\end{longtable}

\subsection{Duplicates and ambiguous names}

Several apparent additions to the frontier are the same literal graph:
five principal cached-controller deliveries repeat PA194C;
PA194X equals PA194CSZ;
the filed PA194HS equals the lowercase packet's PA194NS;
and the two N6 deliveries contain the same PA194N6.
The \emph{uppercase} incoming \file{PA194NS_research.zip} uses PA194NS
for a different machine, preserving the sum guard rather than testing the
cache alone. Its dense digest begins \texttt{eef9329b}; the filed HS/NS
digest begins \texttt{19935635}. HAZ's preliminary H is the old MPH auxiliary,
not the H from the HS packet. Full paths and exact bytes resolve these
collisions; names alone do not.

\subsection{The larger machines}

The 19 graphs at 195 states include the original bundled-diagonal PA195,
the separated-head 195S, the cached 195C, the scanner variant M195S8,
diagonal-combined PA195DC, quotient parents and rejected or auxiliary
experiments. Two latest QN-tail experiments are physically checked but are
not supplied with complete new consistency calibrations. Thus 19 is not
a count of 19 independently established consistency theorems.

The 99 graphs at 196 states include the imported PA196 and later logical,
register, controller and decoder variations. A shorter direct-negation
source has a verified 197-state allocation: source length is not the
physical objective. The 197--207-state graphs mostly record larger
experiments or ancestors. There are 142 classes at 208 states, predominantly
U208 research variants and quoted inputs; the principal U208 reductions use
entry, finite-history and contextual folds of an older U209 implementation.
Imported PA209 and U209 are distinct source lineages. The census retains
all recognized ancestors through 222 states and their exact paths.

\subsection{A useful negative control}

Three distinct comparator-return mutations independently halt after exactly
\textbf{1,361 blank-tape steps}. The fresh simulator and their identities are
included in \file{comparison/verify_negative_examples.py} and
\file{comparison/early_halt_rechecks.json}. Their source-level zero invariants
do not rescue a changed initialization. This is why a successful source
contract and a full blank initialization are both required below. The fourth
rejected 194-state graph is the shortened seven-zero scanner, whose supplied
source-address contract fails; no new complete replay of that older failure
is claimed here.

\subsection{The formal proof boundary}

The current PA project explicitly leaves its machine-specific iPRA
consistency reductions open, while providing substantial Lean source,
quotient, zipper, scheduler and clock infrastructure~\cite{pa-status}.
The delivered papers separately give ordinary mathematical arguments
intended to be formalizable in PA. These are different claims. Absence of
a newly compiled Lean endpoint does not disprove an ordinary paper proof;
standard-natural-number simulation and a conditional arithmetic endpoint
also do not discharge the missing iPRA reductions. This report proves its
new transport and identifies the inherited PA-calibration arguments rather
than silently upgrading one proof level to another.


\section{Logical background and the inherited PA calibration}
\label{sec:logical-background}

The numerical syntax uses Cantor pairing
$p(a,b)=(a+b)(a+b+1)/2+b$ and constructors
\[
 I(A,B)=4p(A,B),\quad M(a,b)=4p(a,b)+1,\quad
 U(a,B)=4p(B,a)+2,\quad E(a,b)=4p(a,b)+3.
\]
Code zero denotes truth, code one denotes the chosen contradiction, and
$N(B)=I(B,1)$. Write $A\to B$ for $I(A,B)$, and use the usual classical
abbreviations for conjunction and existential quantification. The six
universally instantiated heads, in their order of emission, are
\begin{align*}
 h_1(D,A,H)&=(\exists x_D\,U(D,H))\to H,\\
 h_2(D,A,H)&=U(D,U(A,H))\to U(A,U(D,H)),\\
 h_3(D,A,H)&=M(H,D)\to(E(H,A)\to M(A,D)),\\
 h_4(D,A,H)&=E(H,D)\to(M(A,H)\to M(A,D)),\\
 h_5(D,A,H)&=\exists x_D\,E(H,D),\\
 h_6(D,A,H)&=U(D,A)\to\bigl(U(D,A\to H)\to U(D,H)\bigr).
\end{align*}
The Q emitter replaces just the last of these by
\[
 h_6^{\mathrm Q}(D,A,H)
 =U(D,A\to H)\to\bigl(U(D,A)\to U(D,H)\bigr).
\]
All other axiom heads, the diagonal schema
$\Gamma(B)=B\to U(B,B)$, the guarded separation gate and the inference
rules are retained. The propositional head is
\[
 \rho(D,A,H,d)=((H\to d)\to D)\to((D\to H)\to(A\to H)),
 \qquad d\le D.
\]
The guarded gate is $G(a,b,c)=N(V(a,b,c)\to Q(a,b,c))$ for $c\ne a,b$,
with a failed guard producing zero, where
\begin{align*}
 V(a,b,c)&=M(a,b)\to U(c,M(a,b)),\qquad P(a,b,c)=c\to M(c,b),\\
 S(a,b,c)&=U\bigl(b,(M(a,b)\to P(a,b,c))\wedge(M(c,b)\to c)\bigr),\\
 Q(a,b,c)&=U(c,S(a,b,c)\to N(P(a,b,c))).
\end{align*}
The rules are modus ponens and generalization of a theorem
$B\mapsto U(i,B)$, together with truth. Proofs are finite trees with
explicit axiom parameters and inference annotations.

\subsection{Why the exchanged head preserves the calculus}

The unchanged dominated implication basis derives the $K$ and $S$
schemata, hence the uniform exchange theorem
\[
 (P\to(Q\to R))\to(Q\to(P\to R)).
\]
The relevant supplied Lean source is
\file{M6/Logic/RestrictedImplication.lean}; its hypotheses are modus ponens,
the dominated row, reflexivity of the numerical bound and propagation of
bounds into either implication argument~\cite{restricted}. No version of
the sixth head is used to obtain exchange. Instantiating exchange once
derives either sixth head from the other. Recursively replacing the
changed axiom leaves therefore gives primitive-recursive proof translations
in both directions, preserving conclusions literally. PA verifies these
translations by induction on finite derivations.

The new \file{proofs/exchange_certificate.json} is an additional closed
61-node $K,S,\mathrm{MP}$ derivation: 17 $K$ instances, 14 $S$ instances and
30 modus-ponens steps. Its small replay checker validates the published DAG
separately from the bracket-abstraction program that generated it. This
certifies the finite exchange template; the derivation of $K,S$ from the
dominated numerical basis remains the identified inherited theorem.

\subsection{The no-push command interface}

Let $[a;w]=p(a,w)$, with total Cantor projections. In particular
$[0;0]=0$ and both projections of zero are zero. A command store is
$(t,w,d,h)$: current formula, numerical stack, offset and cached formula.
For $y=d+x$, QN and QNS implement
\[
 \mathsf{MP}_x(t,w,d,h)=
 \begin{cases}(y,w,y,0),&I(t,y)=h,\\(0,w,y,0),&I(t,y)\ne h,\end{cases}
\]
and
\[
 \mathsf S(t,w,d,h)=
 \begin{cases}\mathsf{halt},&t+h=1,\\
 (h,\operatorname{tl}(w),d,\operatorname{hd}(w)),&t+h\ne1.
 \end{cases}
\]
The altered physical comparator preserves this entire command, including
its original sum guard. It does not replace it by a cache-only test.

With $D=d+x$ and $H=h+y$, the axiom command has the same ten-entry interface:
\[
 \mathsf{AX}_{A,x,y}(t,w,d,h)
  =\bigl(g(H,D,A),[q_1,\ldots,q_9,t;w],0,0\bigr),
\]
where $q_1=\Gamma(A)$, $q_2=h_6^{\mathrm Q}$, the next five entries are
$h_5,\ldots,h_1$, $q_8=\rho(D,A,H,d)$ and $q_9=U(d,h)$.
Each freshly generated code $g,q_1,\ldots,q_9$ is different from one;
the saved incoming top $t$ is unrestricted. The complete proof compiler
must reserve the zero that successful MP no longer pushes. The exact
all-store Push, Pop, Discard, selector, generalization and MP contracts are
retained in \file{construction/references/logical/pa194n_command_proof.tex}.
They use total projections rather than treating the numerical stack as a
free list, and provide a halting proof word from every incoming store for
each derivation of code one. Cofinal axiom padding makes the unchanged
Cantor scheduler eventually execute such a word.

\subsection{The two inherited syntactic reductions}

The separate-head PA calibration has two distinct directions~\cite{primary,qn}.
For completeness, the diagonal clause recovers the reverse membership
rewriting and neighborhood-equality interfaces used by the original
PA-refutation compiler. This produces a proof of code one in the numerical
calculus from each PA refutation. For soundness, translate membership as
finite binary membership, equality as equality and $U$ as quantification.
PA supplies finite binary comprehension by induction, via a
primitive-recursive \emph{proof-producing} template. Code freshness
justifies the diagonal axiom. The $a=b$ separation case and the zero-code
alias $I(0,0)=0$ are handled separately in those supplied arguments.

Annotate each live theorem code and canonical stack entry with a PA proof
of its translation. The command operations preserve this annotation, and
a halt gives a PA refutation. This uses finite syntax and proof checking;
it does not introduce an arithmetic truth predicate or appeal to PA
reflection. Together these syntactic reductions give the inherited
PA-provable equivalence for QN. The new physical and source transport is
given in full next.

% Report fragment: requires amsmath, amssymb, amsthm, booktabs, longtable.
\section{PA194QNS: the certified fallback construction}

The machine PA194QNS composes the supplied Q axiom emitter, the supplied
no-push modus-ponens interface, and the supplied SUMS implementation of the
original sum-based Shift guard. None of those components is claimed to be
new individually. The new object is their exact 194-control composition,
with complete source and tape certificates and a finite-trace equivalence
to the pinned PA194QN construction. The equivalence is an ordinary
mathematical proof formalizable in PA; no new Lean or Rocq kernel endpoint
is claimed.

\subsection{Literal identity and state budget}

The tape alphabet is $\{0,1\}$, the tape extends in both directions, and
every cell is initially zero. The initial head coordinate is zero. The
working controls are $0,\ldots,193$, with initial control zero and separate
terminal halt label $194$. Every one of the 388 actions writes a bit and
moves exactly one cell left or right. The dense table's SHA-256 is
\begin{quote}\ttfamily\small
7eb36aea69e5de46cc47fa7aab0191f191657ab500dafb808dc3c394640591b8
\end{quote}
The native-label table starts at 74 and halts at 254. Its SHA-256 is
\begin{quote}\ttfamily\small
11613b0b12395fd3132bfe5038c47297685adf337459b138879bb8cdf790daba
\end{quote}
The delivered bijection maps the complete native table to the dense table;
all 194 working controls are graph-reachable.

The unfolded parent has 135 fixed controls and 67 AX readers. The seven
Q-emitter folds and the added comparator fold are simultaneously compatible:
\[
 135+67-8=194.
\]
The 32 local history/read cases are checked against the final simultaneous
quotient. The arithmetic-box rows are unchanged. The separate literal
symbolic checker derives its four list-sweep lemmas from the actual table,
rather than calling the first checker's inherited arithmetic-box oracle.

\subsection{Exact patch relative to QN}

The frozen parent PA194QN has dense digest
\begin{quote}\small\nolinkurl{d1cf0f6e7a3c8f999d045694efc107dd1f9a03d45dd6c15e3d197ff7754c7231}.\end{quote}
The builder reconstructs its entire native and dense tables and source
instruction map before applying the new transformation. The final native
patch is exactly the following; every other row is identical to QN.
\begin{center}
\begin{tabular}{rrll}
\toprule Control&Read&QN action&QNS action\\\midrule
28&0&$(0,R,41)$&$(0,R,330)$\\
38&1&$(1,R,20)$&$(1,R,30)$\\
41&0&$(0,R,24)$&deleted\\
41&1&$(0,R,25)$&deleted\\
70&0&$(0,R,41)$&$(0,R,330)$\\
310&1&$(1,R,62)$&$(1,R,330)$\\
330&0&absent&$(0,L,31)$\\
330&1&absent&$(1,L,46)$\\\bottomrule
\end{tabular}
\end{center}
The no-push action $(164,1)\mapsto(1,R,114)$ is already in QN and is
retained. This avoids claiming it again as a new row change.

Before quotienting, added control 330 has actions $(0,R,63)$ and
$(1,L,73)$. Incoming edges to old control 41 write zero and move right;
incoming edges to raw 330 write one and move right. The immediately
preceding cell therefore distinguishes the two histories. Folding
$41\to330$ gives the displayed final actions and the four three-step
local traces
\[
330,31,21,24;\quad 330,46,61,25;\quad
330,31,47,63;\quad 330,46,62,73.
\]
They cover histories zero/one and current bits zero/one, preserve the
history cell and realize precisely the raw writes, moves and target
controls. The other seven folds retain their full certificates in the
presence of this additional fold.

\subsection{An exact source equivalence}

Both programs use logical registers
$(N,T,W,A,D,H,X,Y,L)$, indexed from zero. Their physical allocation is
$(6,2,0,4,3,7,8,5,1)$. QN has 3,654 primitive source instructions; QNS
has 3,651. Their instructions numbered $0$ through $3648$ are literally
equal, including every resolved target. Only the final five-instruction
tail is replaced by two instructions:
\begin{center}
\begin{tabular}{p{.08\linewidth}p{.43\linewidth}p{.39\linewidth}}
\toprule Site&QN&QNS\\\midrule
3649&decrement $X$; zero $\to3652$, positive $\to3650$&increment $H$; $\to3650$\\
3650&increment $H$; $\to3651$&jump $0$\\
3651&jump $3649$&absent\\
3652&increment $X$; $\to3653$&absent\\
3653&jump $0$&absent\\\bottomrule
\end{tabular}
\end{center}
The only edge from the common prefix into either tail is the zero branch
of instruction 3645, which decrements $H$. Thus $H=0$ on every tail entry
arising from the designated initial source instruction 25. This is checked
by a separate structural certificate that imports neither compiler nor
physical executor.

The common instructions $0,\ldots,10$ form a destructive equality test:
on input $(H,X)=(u,v)$ they halt exactly when $u=v$; otherwise they
terminate at instruction 11 with $H=X=0$ and all other registers
unchanged. Lockstep decrements detect simultaneous exhaustion; the
remaining loops clear an unmatched counter. Induction on $u+v$ proves
termination and the asserted result in PA.

At entry 3649 write $X=x$; the incoming invariant gives $H=0$. QN's
finite tail transfers $X$ to $H$ and then increments $X$, yielding
comparator operands $(x,1)$. QNS simply increments $H$, yielding
$(1,x)$. They therefore both halt exactly when $x=1$, and otherwise
return to the identical complete store at instruction 11. The transfer
loop terminates by induction on its initial value $x$.

More explicitly, the common Shift prefix at 3642 transfers incoming
$T=t$ into $X$, then incoming $H=h$ into both $T$ and $X$. Starting with
arbitrary $(T,H,X)=(t,h,x)$, it reaches the tail with
\[
 (T,H,X)=(h,0,x+t+h).
\]
Both complete Shift blocks consequently halt exactly when $x+t+h=1$;
otherwise they leave $(T,H,X)=(h,0,0)$ and preserve every other register.
At command boundaries $X=0$, recovering precisely the original guard
$t+h=1$. No cached-head-only guard is being substituted.

For finite-trace transport, cut a source execution whenever it enters
3649 and follow that local block through the comparator to instruction 11
or halt. Replace each such terminating block by its counterpart. Outside
these blocks the resolved instruction graphs are identical. The sole-entry
check supplies $H=0$ at every replacement, and every block has positive
finite duration. Induction on the number of replacements gives
primitive-recursive translations of finite halting traces in both
directions, from any complete register store at source 25. PA proves
the correctness of these translations. In particular,
\[
 \mathrm{PA}\vdash
 \operatorname{Halt}(\mathrm{source}_{\rm QNS},25)
 \leftrightarrow
 \operatorname{Halt}(\mathrm{source}_{\rm QN},25).
\]

\subsection{From the source to consistency}

The complete blank computation of QNS reaches the clean source-25 frame
after 58,004 physical steps, at native control 84 and head coordinate
$-2$. The logical registers are $(2,0,0,0,0,0,0,0,0)$, followed by 26
one-valued spare registers. Both simulators compare the entire finite
tape support and zero exterior, not a tape window. QN reaches that
source frame after 58,154 steps.

The universal source-frame certificates cover 4,751 primitive outcomes,
including both branches of every decrement, arbitrary natural primary
registers, an arbitrary affine header gap, and the stated complete spare
list interface. Together with positive progress and deterministic tape
execution, they identify blank halting with halting of the corresponding
source from its initial frame. The new source equivalence therefore
transports the supplied QN paper calibration:
\[
 \mathrm{PA}\vdash
 \operatorname{Halt}(\mathrm{PA194QNS},\mathrm{blank})
 \leftrightarrow \exists e\,
       \operatorname{Prf}_{\mathrm{PA}}(e,0=1),
\]
and hence
\[
 \mathrm{PA}\vdash
 \neg\operatorname{Halt}(\mathrm{PA194QNS},\mathrm{blank})
 \leftrightarrow\operatorname{Con}(\mathrm{PA}).
\]
The inherited QN calibration retains the separate diagonal-head calculus,
the K/S exchange derivation for Q's curried-premise change, and the
reserved-zero proof compiler for no-push MP. They are attributed in the
reference fragments and are not proved again merely by executing the
physical certificates. This transport needs no assumption that PA is
consistent. Such an assumption is used only to conclude actual
nonhalting from the biconditional.

\subsection{Material difference and its limits}

QNS changes the executing Shift protocol of QN, eliminating a counter
transfer on every visit while retaining its arbitrary-store behavior.
It changes the Q axiom emitter relative to the supplied sum-symmetric
NS machine. These are recurring logical and runtime operations, not
unreachable-row edits. Its complete action graph is absent from all 39
current 194-state graph classes at the pinned revision, even after
state renaming, spanning-root change, reflection and bit interchange.
The component techniques are inherited, and no sub-194 construction or
global minimality result is claimed for this fallback.


\section{Fresh verification and what it proves}
\label{sec:verification}

\begin{center}\small
\begin{tabularx}{\linewidth}{@{}p{.50\linewidth}X@{}}\toprule
Check & Fresh result\\\midrule
Working states / transitions & 194 / 388\\
Raw controls / simultaneous folds & 202 / 8\\
Local fold cases & 32\\
Primitive source sites / outcomes & 3,651 / 4,751\\
First symbolic executor & 4,751 outcomes; 1,512,246 symbolic steps\\
Independent literal executor & 4,751 outcomes; 1,106,186 symbolic steps\\
Generic list lemmas independently rederived & 4\\
Minimum positive duration constant & 3\\
Blank initialization & 58,004 physical steps; entire tape equal\\
Prior 194-state classes excluded by exact graph comparison & 39\\
Source prefix identical to QN & 3,649 instructions\\
Exchange proof DAG & 61 nodes; no hypotheses\\
\bottomrule\end{tabularx}
\end{center}


The two symbolic checks have deliberately different trusted interfaces.
The first executes the exact source contracts using inherited arithmetic
block lemmas. The second imports no allocator or source compiler: it
reconstructs the control-flow graph twice, proves four generic list sweeps
from the literal rows, and executes every primitive outcome on symbolic
complete tapes. Both quantify over arbitrary natural register values;
they are not bounded tests on small register contents. The independent
source-tail certificate checks the actual comparator rows, sole-entry
invariant and identical source prefix without importing either physical
executor.

The blank-start checks are finite, but they compare the whole finite tape
support, head and control, with a zero exterior. The universal source
simulation then applies from that exact seed. Every simulated source step
takes positive time. Determinism and this positive progress give the
converse halting implication as well: a first physical halt cannot be
passed by arbitrarily many positive simulated nonhalting steps.

These are executable mathematical certificates plus the stated ordinary
proof. The JSON receipts explicitly record \file{kernel_checked=false}.
No test is described as a completed Lean/Rocq proof of the PA or iPRA
consistency equivalence, and no bounded nonhalting run is used to establish
one.

\section{Reduction research and a different numerical encoding}
\label{sec:research}

\subsection{What was searched below 194}

The first bounded comparison combines two AX grammars (primary and Q),
four constructor coefficients (4, 5, 6, 8), four comparator/controller
layouts and two final-return choices: 64 explicit combinations. For each,
the program enumerates its one-bit incoming-history fold opportunities and
checks compatible disjoint subsets against the final simultaneous table.
The best physical count is 194. These are allocation screens, not blanket
correctness certificates for every proposed command combination. Only the
separately accepted constructions receive their complete interfaces and
physical checks.

The SUMS and coefficient-5/6 combinations expose a concrete obstruction to
adding previously published savings: both desired folds remove native
control 41. Two new helpers cannot be paid for by deleting that one control
twice. Their naive composition costs 195. A coefficient-eight construction
also remains at 195 in the screened layout.

The new QNS table was additionally subjected to three explicit local
contraction searches with radius six and a 256-step detour limit. The
forward local abstraction contains 275,931 state/window contexts, starting
from all 13-bit windows at native start 74. The searches find no accepted
single same-head merge, no accepted two-square head-displacement merge,
and no accepted adjacent-marker-bit merge. The latter two screens examine
32,367 and 32,173 eligible ordered pair/offset cases, with 1,565,312 and
914,435 action trials respectively. Complete output and the executable
JavaScript are included.

These local abstractions overapproximate the encountered tape windows;
the search tests a sufficient, very specific equivalence pattern. A negative
result says nothing about more global tape recodings, changes to several
helpers at once, a larger radius or a different proof calculus. It must not
be described as a lower bound for arbitrary PA-consistency machines.

\subsection{Exchange of other logical premises}

The Q transformation exchanges the two premises of quantifier distribution.
The same uniform exchange proof also permits
\begin{align*}
 h_3^{\mathrm e}(D,A,H)
  &=E(H,A)\to\bigl(M(H,D)\to M(A,D)\bigr),\\
 h_4^{\mathrm e}(D,A,H)
  &=M(A,H)\to\bigl(E(H,D)\to M(A,D)\bigr).
\end{align*}
These change the membership-rewriting heads, rather than merely renaming
their universally quantified parameters. Replacing any subset of
$h_3,h_4,h_6$ by its exchanged form preserves numerical derivability by
primitive-recursive translations in both directions. The independent
free-algebra audit verifies all ten pushed AX entries and the exact guarded
top for each proposed emitter. The 190-token words were then searched for
different associations, repeated subtrees, compatible folds and parameter
permutations. The final 330-second run made 1,727,732 proposal slots and accepted
342,713; its best projected allocation was $134+71-8=197$. No candidate
at or below 194 was found in that run. The proof of
interderivability is useful independently of the allocation result.

\subsection{TAG196N: a new formula-code involution}

The secondary delivered construction changes the numerical syntax to
\[
 I'(A,B)=4p(A,B),\quad M'(a,b)=4p(a,b)+1,\quad
 U'(a,B)=4p(B,a)+3,\quad E'(a,b)=4p(a,b)+2.
\]
Only the quantifier and equality tags are exchanged; truth zero and
contradiction one are fixed. This is a genuine change of the numbers
constructed and compared during proof search. The stack and command
scheduler retain their old pairing representation.

Let $\sigma$ exchange 2 and 3 and fix 0 and 1. Define $F(0)=0$, and for
$n>0$ write uniquely $n=4p(a,b)+r$, $0\le r<4$, and put
\[
 F(n)=4p(F(a),F(b))+\sigma(r).
\]
Both children are below $n$, so this is primitive-recursive
course-of-values recursion. Strong induction gives the four constructor
identities, with the \emph{same} $F$ applied to variable labels and formula
codes. A second induction gives $F(F(n))=n$. Thus $F$ is a bijection of all
natural codes; it fixes zero and one and preserves the distinct-label gate.
In particular, the binder label in the recoded diagonal instance is
$F(B)$, exactly the code of its recoded property. Recoding formulas alone
would not prove that identity.

The translation is not order preserving: $F(2)=3$ and $F(3)=2$. Therefore
one may not transport the side condition $d\le D$ by monotonicity. At a
dominated-row leaf, use instead the supplied finite derivation of the
\emph{unrestricted} $\rho$ instance. The component bounds for
$I(A,B)=4p(A,B)$ remain valid in the target. This supplies proof translations
in both directions without a false numerical-order assumption. Translating
a derivation and then compiling a fresh target command word is essential;
an old word is not simply asserted to have the same execution.

The separate parameter screen completed 139,968 assignments. The final
330-second tag/Gamma search made 1,647,198 proposal slots and accepted
268,038; its best allocation remained 196. Proposal slots can repeat a
grammar or make no change, so these counters are not counts of distinct
machines. The accepted physical construction is verified independently of
its search score.

The changed AX word has 202 tokens: its 15 quantifier tags each gain an
increment, while its three equality tags each lose an increment. Exposing
the existing triple-increment reader as a reusable AX macro is a useful
allocation improvement. The delivered TAG196N budget is
\[
 135\text{ fixed controls}+71\text{ AX readers}
      -10\text{ simultaneous folds}=196.
\]
It uses no-push MP and the original sum guard. Its source has 3,666 sites;
all 4,767 universal primitive outcomes pass the first symbolic executor,
with 1,535,770 symbolic operations and maximum 517. Complete blank
initialization takes 58,320 steps. The arithmetic-box rows are unchanged
dependencies of this check; the full second literal-checker run performed
for QNS is not claimed for TAG196N. A separate no-compiler audit checks
the entire 19-instruction no-push splice, and a free-constructor audit
checks every emitted formula and the ten-entry stack interface. A further
source-binding audit checks all 46 arithmetic-box rows against the frozen
parent, all 61 constructor prefixes and four-increment bodies, and equality
of the complete resolved primitive program before and after reassociation.

Its complete tables, deterministic builder, splice and emitter checks,
mathematical recoding and command-compiler proofs, receipts and licenses
are under \file{research/tag196n/}. The dense digest is
\begin{quote}\small\nolinkurl{e250b15fb2ad284cffaab5ed2842593f51d2ad7eec6c707aa42786c5ce3cfac6}.\end{quote}
The table is distinct from all 99 existing 196-state action graphs even
under the stronger cosmetic comparison. This secondary result changes
the arithmetic encoding, but \textbf{does not improve the 194-state bound}.
The archive includes its universal mathematical proof separately from the
bounded recoding calculations; those calculations are supporting audits,
not proofs of involution for all natural numbers.

\subsection{What transfers from the ZFC and ConciseIPRA work}

The reusable ZFC ideas are exact implementation mechanisms: caller-history
folding, simultaneous allocation checks, command-local field readers,
compiler-supplied spare zeros and hereditary recoding of formulas
\emph{and} labels~\cite{zfc}. The no-push protocol here has that documented
ancestry. The exact K215 external Lean/Rocq equivalence and its internal PA
endpoint have different statuses; a proof for the former is not silently
treated as a proof of the latter. Likewise, a ZFC command saving can conflict
with the actual PA reader contexts and cannot simply be subtracted from
the PA count.

ConciseIPRA shows how a small native calculus can derive familiar logic and
how to remove redundant certificate machinery~\cite{ipra}. Its objective
is the size of a complete Python proof checker, and its calibration target
is iPRA. Both differ from a binary one-tape implementation searching PA
refutations. Built-in substitution, binding and list operations still cost
machine controls, and replacing the PA-strength proof supplier requires a
new PA refutation translation. The useful lesson is to optimize the
calculus, certificate interface and literal runtime together; no direct
state saving or consistency-strength equivalence is inferred from a short
Python checker.


\section{Further research questions}

\begin{enumerate}
\item Can a decoder grammar and its compatible contextual folds be
optimized jointly, including fixed arithmetic readers? Minimizing only the
raw grammar can worsen the final physical budget. The new code and axiom
experiments provide additional exact search domains.
\item Can a larger local relation, or a simultaneous change of multiple
helpers, remove a state that the radius-six, 256-step single-contraction
model cannot remove? The present negative result only concerns that
sufficient local model.
\item Can the unused zero-push work be traded for a different physical
register/continuation interface that actually removes a control? Deleting
source instructions alone has repeatedly left the state budget unchanged.
\item Can the formula-tag involution be allocated in 194 or fewer states,
while retaining a checked command compiler? Its semantics are particularly
simple, but the shared three-increment reader and history-fold constraints
must be priced together.
\item Can the generic PA proof translations and literal simulation be
internalized in the repository's intended iPRA endpoint? A useful next
formal deliverable is the short QN--QNS tail equivalence, followed by the
already identified common calibration obligations.
\end{enumerate}

\section{Reproduction and archive guide}

Run \file{python3 verify_all.py} from the extracted archive. This reconstructs
the principal table in a fresh temporary directory, replays its source and
literal-tape certificates, checks its identity and graph novelty, and
replays the finite exchange proof. Python 3.10 or later and its standard
library are sufficient for the principal checks; assertions must remain
enabled. Search replays additionally use Node.js or a C++ compiler as
specified in \file{search/README.txt}.

The complete table is in \file{construction/machine/}; its source,
instruction map, labels, local fold cases and fresh receipts are adjacent.
\file{comparison/} contains all 54 packet records, all 39 frontier records,
the full census and the 721 table byte presentations needed to audit it.
\file{construction/references/} preserves the attributed mathematical
dependencies, while \file{proofs/} contains the new exchange certificate and
additional proof material. \file{search/} distinguishes accepted outputs
from unaccepted candidates and bounded negative results.

Compile this report with \file{python3 build_report.py}, or repeat
\file{pdflatex} on \file{report/pa194_qns_research.tex} from its report
directory until cross-references stabilize. The source files, both final PDFs, manifests and licenses are included.
The supplementary \file{report/calculus_proofs.pdf} presents the full
logical transports and the matching no-push command compiler.
The archive is self-contained for the main construction and checks; only a
fresh full repository census requires access to the pinned repository.

\appendix
\clearpage
\section{Every existing 194-state graph}
\label{app:frontier}
This is the pre-construction frontier at the pinned snapshot. Exact hashes,
representative paths and every occurrence are in the accompanying CSV/JSON
records. ``Reported construction'' describes the paper's evidence level;
it does not assert completed iPRA formalization.
\begingroup\small
\begin{longtable}{@{}p{.23\linewidth}p{.14\linewidth}p{.57\linewidth}@{}}
\toprule Name & Evidence status & Distinguishing construction\\\midrule\endhead
CS both-zero6 & auxiliary & Known-zero X decrement at both source26 and final emitter exit; research composition.\\
CSZ-A3 & auxiliary & CSZ plus comparator-return register3 and final-emitter register3 decrements.\\
CSZ-A5 & auxiliary & CSZ plus comparator-return register3 and final-emitter register5 decrements.\\
CSZ-A6 & auxiliary & CSZ plus comparator-return register3 and final-emitter register6 decrements.\\
Cached I-tree zero-jump seed & auxiliary & Alternate decoder plus cached controller and entry zero jump; search basis.\\
Gamma(D)-CSZA & auxiliary & Supplementary Gamma(D) emitter plus known-zero changes.\\
MPH & auxiliary & MP result cached in H while retaining original Shift; physical search input.\\
PA194CZ & auxiliary & Cached controller; entry known-zero decrement of register 6.\\
PA194altC & auxiliary & Cached controller on alternate 68-reader decoder.\\
Rejected CS node1-X & rejected & Native103/read1 becomes dec6; exact blank run halts at1361.\\
Rejected I-tree node1-X & rejected & Known-zero source rewrite fails blank bootstrap; exact halt at1361.\\
Rejected ZR node1-X & rejected & Register6 comparator-return variant; exact blank run halts at1361.\\
Rejected seven-zero scanner & rejected & Decoder local folds pass, but source-address scanner violates runtime contract.\\
M4R & reported construction & Reassociated coefficient4 constructor decoder; original resolved source graph.\\
M5 & reported construction & Formula constructor coefficient5; explicit variable and proof recoding.\\
M6 & reported construction & Formula constructor coefficient6; explicit variable and proof recoding.\\
PA194 & reported construction & Separate diagonal head, 67-reader decoder with seven folds, original controller.\\
PA194-68/8 & reported construction & Same separate-head source, reassociated 68-reader decoder with eight folds.\\
PA194-AD & reported construction & PA194 with physical A/D register exchange.\\
PA194A & reported construction & CS entry jump at source26 becomes decrement of known-zero register A.\\
PA194C & reported construction & MP result remains cached in H; Shift tests H=1; original ten-entry logical interface recompiled.\\
PA194CS & reported construction & Symmetric Shift guard; cached MP controller; contextual controller fold.\\
PA194CSZ = PA194X & reported construction & CS entry jump at source26 becomes decrement X; two headline names denote identical dense bytes.\\
PA194D & reported construction & Separate diagonal axiom uses Gamma(D), new proof selector, 69 readers and nine folds.\\
PA194F & reported construction & CS final emitter jump becomes decrement of known-zero register X.\\
PA194H (HS packet) & reported construction & No-push successful MP with cached-head-only Shift; asymmetric comparator.\\
PA194HA & reported construction & MP result cached in H with original sum guard, plus physical A/D placement exchange.\\
PA194HAZ & reported construction & HA plus known-zero final AX decrement; full startup and separate four-site source invariant.\\
PA194HS = PA194NS (cached) & reported construction & No-push successful MP with cached-head-only Shift; symmetric comparator. HS and lowercase NS packets give identical dense bytes.\\
PA194J & reported construction & Ten source jumps become known-zero decrements; three physical rows change.\\
PA194N & reported construction & Successful MP drops dummy-zero push, retains original Shift; revised command proof compiler.\\
PA194N5 & reported construction & Coefficient-five formula constructors combined with no-push MP and original sum guard.\\
PA194N6 & reported construction & Coefficient-six constructors combined with no-push MP; joint N5/N6 and standalone N6 packets give identical dense bytes.\\
PA194NS (sum) & reported construction & SUMS comparator with no-push MP. This is a different graph from the cached-only NS/HS machine despite the name collision.\\
PA194Q & reported construction & Quantifier-distribution curried premises exchanged; K/S proof translation.\\
PA194QN & reported construction & Exchanged curried quantifier-distribution premises combined with no-push MP and original sum guard.\\
PA194R & reported construction & Reassociated decoder of PA194F; 1111 source paths changed, same resolved source graph.\\
PA194SUMS & reported construction & Original sum guard with operands of destructive comparator reversed; one new history fold absorbs replacement trailer.\\
PA194ZR & reported construction & Eight comparator returns and final emitter return become known-zero register3 decrements.\\
\bottomrule\end{longtable}

\endgroup

\clearpage
\section{Complete principal transition table}
\label{app:table}
For each state the two actions correspond to reading zero and one. An
action $w\,D\,q'$ writes $w$, moves $D\in\{L,R\}$ and enters $q'$.
State zero is initial and 194 is terminal. The two side-by-side state
ranges are only a typesetting convention; together they give all 388 rows.
\begingroup\small
\renewcommand{\arraystretch}{1.0}
\begin{longtable}{@{}rcc@{\hspace{2.5em}}rcc@{}}
\toprule State & Read 0 & Read 1 & State & Read 0 & Read 1\\\midrule\endhead
0 & $1\,R\,1$ & $1\,R\,2$ & 97 & $0\,L\,45$ & $0\,R\,108$\\
1 & $1\,R\,3$ & $1\,R\,1$ & 98 & $0\,L\,44$ & $1\,R\,113$\\
2 & $0\,R\,0$ & $1\,R\,2$ & 99 & $0\,L\,114$ & $1\,R\,115$\\
3 & $0\,R\,4$ & $0\,R\,1$ & 100 & $0\,R\,116$ & $0\,L\,14$\\
4 & $1\,R\,5$ & $1\,L\,6$ & 101 & $0\,R\,109$ & $1\,R\,101$\\
5 & $1\,L\,7$ & $1\,L\,8$ & 102 & $0\,L\,45$ & $0\,R\,101$\\
6 & $0\,L\,7$ & $1\,L\,6$ & 103 & $0\,L\,114$ & $1\,L\,33$\\
7 & $0\,L\,9$ & $1\,L\,6$ & 104 & $0\,R\,66$ & $0\,R\,117$\\
8 & $0\,L\,10$ & $0\,L\,11$ & 105 & $0\,R\,56$ & $0\,R\,117$\\
9 & $1\,R\,12$ & $0\,L\,9$ & 106 & $0\,L\,62$ & $1\,R\,118$\\
10 & $0\,L\,9$ & $0\,L\,10$ & 107 & $0\,R\,119$ & $1\,R\,120$\\
11 & $1\,L\,8$ & $1\,L\,11$ & 108 & $0\,R\,14$ & $1\,R\,108$\\
12 & $0\,R\,12$ & $0\,R\,13$ & 109 & $0\,R\,117$ & $1\,R\,109$\\
13 & $0\,L\,14$ & $0\,R\,13$ & 110 & $0\,R\,71$ & $1\,R\,121$\\
14 & $1\,L\,15$ & $0\,L\,15$ & 111 & $0\,R\,122$ & $1\,L\,123$\\
15 & $0\,L\,15$ & $1\,L\,16$ & 112 & $0\,R\,124$ & $1\,R\,125$\\
16 & $0\,L\,17$ & $1\,L\,16$ & 113 & $0\,R\,47$ & $0\,L\,14$\\
17 & $0\,L\,18$ & $1\,L\,16$ & 114 & $0\,R\,90$ & $1\,R\,98$\\
18 & $0\,L\,19$ & $1\,L\,16$ & 115 & $0\,R\,114$ & $1\,L\,93$\\
19 & $0\,L\,20$ & $1\,L\,16$ & 116 & $0\,R\,126$ & $1\,R\,126$\\
20 & $0\,L\,21$ & $1\,L\,16$ & 117 & $0\,R\,127$ & $1\,R\,117$\\
21 & $0\,L\,22$ & $1\,L\,16$ & 118 & $0\,R\,128$ & $1\,R\,129$\\
22 & $0\,L\,23$ & $1\,L\,16$ & 119 & $0\,R\,32$ & $1\,R\,64$\\
23 & $0\,R\,23$ & $1\,R\,24$ & 120 & $0\,R\,31$ & $1\,R\,96$\\
24 & $0\,R\,25$ & $1\,R\,26$ & 121 & $0\,R\,48$ & $1\,R\,87$\\
25 & $0\,R\,27$ & $1\,R\,28$ & 122 & $0\,R\,56$ & $0\,R\,127$\\
26 & $0\,R\,29$ & $1\,R\,30$ & 123 & $0\,R\,65$ & $1\,R\,33$\\
27 & $0\,R\,31$ & $1\,R\,194$ & 124 & $0\,R\,97$ & $1\,R\,113$\\
28 & $0\,L\,32$ & $1\,R\,33$ & 125 & $0\,R\,78$ & $1\,R\,77$\\
29 & $0\,R\,34$ & $1\,R\,35$ & 126 & $0\,R\,47$ & $1\,R\,47$\\
30 & $0\,L\,33$ & $1\,L\,36$ & 127 & $0\,R\,67$ & $1\,R\,127$\\
31 & $0\,R\,37$ & $1\,R\,38$ & 128 & $0\,R\,130$ & $1\,R\,131$\\
32 & $0\,R\,39$ & $1\,R\,40$ & 129 & $0\,R\,132$ & $1\,R\,133$\\
33 & $0\,R\,41$ & $1\,R\,42$ & 130 & $0\,R\,134$ & $1\,R\,135$\\
34 & $0\,R\,43$ & $1\,R\,44$ & 131 & $0\,R\,136$ & $1\,R\,137$\\
35 & $0\,R\,45$ & $1\,R\,46$ & 132 & $0\,R\,138$ & $1\,R\,86$\\
36 & $0\,R\,47$ & $1\,R\,48$ & 133 & $0\,R\,139$ & $0\,L\,10$\\
37 & $0\,R\,49$ & $1\,R\,41$ & 134 & $0\,R\,95$ & $1\,R\,123$\\
38 & $0\,R\,50$ & $1\,R\,51$ & 135 & $0\,R\,140$ & $1\,R\,126$\\
39 & $0\,R\,52$ & $1\,R\,53$ & 136 & $0\,R\,141$ & $1\,R\,142$\\
40 & $0\,R\,54$ & $1\,R\,55$ & 137 & $0\,R\,143$ & $1\,R\,144$\\
41 & $0\,R\,56$ & $0\,R\,57$ & 138 & $0\,R\,145$ & $1\,R\,40$\\
42 & $0\,R\,58$ & $0\,L\,14$ & 139 & $0\,R\,146$ & $1\,R\,147$\\
43 & $0\,R\,59$ & $1\,R\,60$ & 140 & $0\,L\,148$ & $1\,L\,143$\\
44 & $0\,R\,52$ & $1\,R\,61$ & 141 & $0\,R\,28$ & $1\,R\,149$\\
45 & $0\,L\,62$ & $1\,R\,63$ & 142 & $0\,R\,135$ & $1\,R\,150$\\
46 & $0\,R\,64$ & $1\,R\,65$ & 143 & $0\,R\,151$ & $1\,R\,152$\\
47 & $0\,R\,66$ & $0\,R\,67$ & 144 & $0\,R\,153$ & $1\,R\,154$\\
48 & $0\,L\,64$ & $1\,L\,68$ & 145 & $0\,L\,130$ & $1\,R\,47$\\
49 & $0\,R\,69$ & $1\,R\,51$ & 146 & $0\,R\,31$ & $1\,R\,50$\\
50 & $0\,L\,70$ & $0\,L\,15$ & 147 & $0\,R\,155$ & $1\,R\,146$\\
51 & $0\,R\,71$ & $0\,L\,10$ & 148 & $0\,R\,156$ & $1\,R\,157$\\
52 & $0\,R\,72$ & $1\,R\,73$ & 149 & $0\,R\,158$ & $1\,R\,159$\\
53 & $0\,R\,74$ & $1\,R\,75$ & 150 & $0\,L\,159$ & $1\,R\,160$\\
54 & $0\,R\,76$ & $1\,R\,77$ & 151 & $0\,R\,40$ & $1\,L\,161$\\
55 & $0\,R\,78$ & $1\,R\,79$ & 152 & $0\,R\,162$ & $1\,R\,163$\\
56 & $0\,L\,80$ & $0\,R\,1$ & 153 & $0\,R\,164$ & $1\,R\,164$\\
57 & $0\,R\,81$ & $1\,R\,57$ & 154 & $0\,R\,165$ & $1\,R\,166$\\
58 & $0\,L\,72$ & $1\,L\,65$ & 155 & $0\,R\,121$ & $1\,R\,111$\\
59 & $0\,R\,82$ & $0\,L\,7$ & 156 & $0\,R\,75$ & $1\,R\,165$\\
60 & $0\,R\,83$ & $1\,R\,84$ & 157 & $0\,R\,167$ & $1\,R\,165$\\
61 & $0\,R\,41$ & $1\,R\,85$ & 158 & $0\,R\,168$ & $1\,R\,169$\\
62 & $0\,R\,86$ & $1\,R\,87$ & 159 & $0\,R\,170$ & $1\,R\,157$\\
63 & $0\,R\,88$ & $0\,L\,71$ & 160 & $0\,R\,135$ & $1\,R\,130$\\
64 & $0\,R\,30$ & $1\,R\,87$ & 161 & $0\,R\,138$ & $1\,R\,135$\\
65 & $0\,R\,69$ & $1\,R\,36$ & 162 & $0\,R\,171$ & $1\,R\,172$\\
66 & $0\,L\,80$ & $0\,R\,2$ & 163 & $0\,R\,173$ & $1\,R\,172$\\
67 & $0\,R\,89$ & $1\,R\,67$ & 164 & $0\,L\,174$ & $1\,R\,175$\\
68 & $0\,R\,90$ & $1\,R\,91$ & 165 & $0\,L\,162$ & $1\,L\,176$\\
69 & $0\,R\,56$ & $0\,R\,81$ & 166 & $0\,R\,161$ & $1\,R\,176$\\
70 & $0\,R\,65$ & $1\,R\,92$ & 167 & $0\,R\,95$ & $1\,R\,177$\\
71 & $0\,R\,41$ & $0\,L\,13$ & 168 & $0\,R\,134$ & $1\,R\,178$\\
72 & $0\,R\,82$ & $1\,R\,68$ & 169 & $0\,R\,150$ & $1\,R\,163$\\
73 & $0\,R\,48$ & $1\,R\,93$ & 170 & $0\,R\,179$ & $1\,R\,180$\\
74 & $0\,R\,41$ & $1\,R\,94$ & 171 & $0\,R\,167$ & $1\,R\,140$\\
75 & $0\,R\,95$ & $1\,R\,96$ & 172 & $0\,R\,134$ & $1\,R\,145$\\
76 & $0\,R\,97$ & $1\,L\,98$ & 173 & $0\,R\,181$ & $1\,R\,40$\\
77 & $0\,R\,48$ & $1\,R\,76$ & 174 & $0\,R\,175$ & $1\,R\,182$\\
78 & $0\,R\,92$ & $1\,R\,99$ & 175 & $0\,R\,167$ & $1\,R\,145$\\
79 & $0\,R\,48$ & $1\,R\,100$ & 176 & $0\,R\,183$ & $1\,R\,173$\\
80 & $1\,R\,80$ & $1\,L\,15$ & 177 & $0\,R\,86$ & $1\,R\,61$\\
81 & $0\,R\,101$ & $1\,R\,81$ & 178 & $0\,R\,184$ & $1\,R\,165$\\
82 & $0\,L\,45$ & $0\,R\,89$ & 179 & $0\,R\,134$ & $1\,R\,181$\\
83 & $0\,R\,102$ & $1\,R\,93$ & 180 & $0\,R\,185$ & $1\,R\,185$\\
84 & $0\,R\,102$ & $1\,R\,103$ & 181 & $0\,R\,167$ & $1\,R\,151$\\
85 & $0\,R\,104$ & $0\,L\,14$ & 182 & $0\,R\,151$ & $1\,R\,130$\\
86 & $0\,R\,105$ & $1\,L\,70$ & 183 & $0\,R\,164$ & $1\,R\,186$\\
87 & $0\,R\,103$ & $0\,L\,14$ & 184 & $0\,R\,140$ & $1\,R\,187$\\
88 & $0\,R\,106$ & $1\,R\,107$ & 185 & $0\,R\,188$ & $1\,L\,189$\\
89 & $0\,R\,108$ & $1\,R\,89$ & 186 & $0\,R\,190$ & $1\,R\,40$\\
90 & $0\,R\,66$ & $0\,R\,109$ & 187 & $0\,R\,32$ & $1\,R\,189$\\
91 & $0\,R\,110$ & $0\,L\,87$ & 188 & $0\,R\,47$ & $1\,R\,191$\\
92 & $0\,R\,30$ & $1\,R\,58$ & 189 & $0\,R\,148$ & $1\,R\,192$\\
93 & $0\,R\,82$ & $1\,L\,111$ & 190 & $0\,R\,40$ & $1\,R\,138$\\
94 & $0\,R\,99$ & $0\,L\,14$ & 191 & $0\,R\,174$ & $1\,R\,47$\\
95 & $0\,R\,64$ & $1\,R\,112$ & 192 & $0\,R\,182$ & $1\,R\,193$\\
96 & $0\,R\,111$ & $1\,R\,74$ & 193 & $0\,R\,161$ & $1\,R\,138$\\
\bottomrule\end{longtable}
\endgroup


\clearpage
\begin{thebibliography}{9}
\bibitem{pa-status} ProveIt, PA consistency machines, README at the pinned
revision. \url{https://github.com/VladimirReshetnikov/ProveIt/blob/317b07268681b44ff352d173a162e2a8c241b2fe/Computability/BusyBeaver/Research/PA/README.md}.
\bibitem{primary} Primary PA194 packet, construction and calibration.
\url{https://github.com/VladimirReshetnikov/ProveIt/tree/317b07268681b44ff352d173a162e2a8c241b2fe/Computability/BusyBeaver/Research/PA/docs/pa194_conpa_2026-11-13}.
\bibitem{qn} PA194QN packet, exact input table, Q emitter and no-push
command calibration. \url{https://github.com/VladimirReshetnikov/ProveIt/blob/317b07268681b44ff352d173a162e2a8c241b2fe/docs/incoming/pa194qn_research_2026-11-13.zip}.
\bibitem{sums} SUMS and sum-guard NS packet, symmetric comparator and
history fold. \url{https://github.com/VladimirReshetnikov/ProveIt/blob/317b07268681b44ff352d173a162e2a8c241b2fe/docs/incoming/PA194NS_research.zip}.
\bibitem{restricted} James Chang, imported PA196 source,
\file{M6/Logic/RestrictedImplication.lean}. \url{https://github.com/VladimirReshetnikov/ProveIt/blob/317b07268681b44ff352d173a162e2a8c241b2fe/Computability/BusyBeaver/Research/PA/jamesyc-upstream/pa196/M6/Logic/RestrictedImplication.lean}.
\bibitem{zfc} ProveIt, K215 source composition and correctness argument.
\url{https://github.com/VladimirReshetnikov/ProveIt/tree/317b07268681b44ff352d173a162e2a8c241b2fe/Computability/BusyBeaver/Research/ZFC/candidate-k215}.
\bibitem{ipra} ProveIt, ConciseIPRA native scalar core and checker-size
objective. \url{https://github.com/VladimirReshetnikov/ProveIt/tree/317b07268681b44ff352d173a162e2a8c241b2fe/Logic/PrimitiveRecursiveArithmetic/ConciseIPRA}.
\bibitem{radix} PA194 constructor-coefficient family and its recoding
calibration. \url{https://github.com/VladimirReshetnikov/ProveIt/tree/317b07268681b44ff352d173a162e2a8c241b2fe/Computability/BusyBeaver/Research/PA/docs/pa194_radix_family_2026-11-13}.
\end{thebibliography}
\end{document}
